\subsection{The exact extension from free terms to flat terms}

\begin{proposition}[A-L18: Proved; the bounded-below coefficient-flat version]
\label{AL18}
In A-L15 replace the assumption that each cochain term
$M^i$ is graded free by the following assumption:
$M^i$ is flat as an ungraded $S$-module and its weight
grading is bounded below. Retain the actual coefficient
inclusion and retraction, the nonnegative cochain
grading, and the stated global cohomology through $d_0$.
Then every conclusion (AL.80) and the residue-field
direct-summand conclusion of A-L15 still hold.
The conclusions (AL.85)--(AL.86) of A-L16 also hold
for every $S$-module $N$, with the same actual
coefficient maps and quotient
$C=M^{d_0+1}/B^{d_0+1}$.
\end{proposition}

\begin{proof}
We first identify exactly the known graded input.
For a nonnegatively graded algebra $D$ over a field
$k$ with $D^0=k$, a bounded-below graded module $L$
with $\operatorname{Tor}^D_1(k,L)=0$ is graded free.
This is the unnumbered remark immediately after
Lemma~2.1 in Herscovich--Rey, Section~2, printed
page~4 of the opened arXiv v1 source
\cite{B-Herscovich-Rey-2012}. Its standing hypotheses
in that section impose no finite-generation condition
on $D$ or $L$.

The actual free presentation used in this known
criterion also makes its application here transparent.
Choose a homogeneous $k$-basis of $L/D^{>0}L$ and
lift its elements to $L$. These lifts define a graded
map
\[
 F=\bigoplus_\lambda D\langle a_\lambda\rangle
                     \longrightarrow L,
 \tag{AL.95}
\]
whose generator weights share the lower bound of $L$.
Induction on weight proves its surjectivity: a
homogeneous element, minus a linear combination of
the chosen lifts, is a finite sum of positive-weight
coefficients times elements of smaller weight.
The induction starts at the common lower bound.
For $K=\ker(F\to L)$, the Tor exact sequence gives
\[
 0\longrightarrow K/D^{>0}K\longrightarrow
 F/D^{>0}F\xrightarrow{\cong}L/D^{>0}L
 \longrightarrow0.
 \tag{AL.96}
\]
The first zero is exactly the specified Tor vanishing.
Thus $K/D^{>0}K=0$, and the same induction on weight
gives $K=0$. This verifies the criterion with the
unrestricted multiplicities and the actual grading
used here. In particular every bounded-below graded
flat $D$-module is graded free.

Take first $R=\mathbb Z_{(p)}$. Since $S$ is free
over $R$, an $S$-flat module is $R$-flat: tensoring
an exact sequence over $R$ first with $S$ and then
over $S$ with that module proves exactness. Consequently
all $M^i$ are $R$-torsionfree, as are their submodules
$Z^i$ and $B^i$. Base change of flatness shows that
$M^i/pM^i$ is flat over $D=S/pS$; its weights remain
bounded below. Applying (AL.95)--(AL.96) makes it a
graded free $D$-module with bounded-below generator
weights. A-L14 therefore makes it graded injective.

The given coefficient splitting makes $B^1$ a
direct summand of $M^0$, so $B^1/pB^1$ is a graded
direct summand of the just constructed free module
$M^0/pM^0$. The exact sequences (AL.81) and their
reductions modulo $p$ are unchanged: their quotients
$B^{i+1}$ are $R$-torsionfree. The identical split
induction used in A-L15 now shows that each
$B^i/pB^i$, $1\leq i\leq d_0+1$, is a graded
direct summand of a bounded-below graded free module,
and is flat after forgetting weights.

Rationally, $M^i_{\mathbb Q}$ is flat over
$\mathbb Q[T]$, hence torsionfree, so its submodule
$B^i_{\mathbb Q}$ is torsionfree and flat over that
principal ideal domain. Each $B^i$ is thus
$R$-torsionfree, with flat reduction modulo $p$ and
flat rationalization. The explicit free-resolution
and torsion-devissage argument (AL.82)--(AL.83),
which was proved for any module with precisely
these three properties, applies without a freeness
assumption on $M^i$. It proves flatness of every
$B^i$ in the required range. As before $Z^0=S$ and
$Z^i=B^i$ for $1\leq i\leq d_0$.

For $R=\mathbb Z$, localization at each rational
prime preserves the given flat modules, lower
weight bounds, maps and cohomological exactness.
The preceding proof applies there, and the actual
Tor-localization argument in A-L15 proves integral
flatness of the boundaries. This proves all the
claims of (AL.80) and the residue-field statement.

Finally the coefficient splitting (AL.87), the exact
sequences with flat boundary quotients, and the
kernel description (AL.88) remain valid. The only
additional property of the middle module used in
(AL.89) is
$\operatorname{Tor}^S_1(M^{d_0+1},N)=0$, supplied
here by its stated flatness. Its Tor exact sequence
therefore gives the identical groups and actual
coefficient splitting (AL.85)--(AL.86).
\end{proof}

This extends the exact module calculation to flat
coefficient terms. An admissible replacement with
such terms and the required strict comparison maps
for every object of the main reduced category is
still an unconstructed step.
